\section{Mechanism of Forking}
\label{sec:mechanism}

The approximate power law in Figure~\ref{fig:frequency} is our first clue: why should context frequency organize the gap so simply?

\subsection{From the frequency law to novel continuations}
\label{sec:coverage}
A rare context has had few chances to reveal its possible continuations. Call a validation target \emph{novel} when its context occurs in training but was never followed there by that target, and let $f(F)$ be the fraction of validation positions at context count $F\ge1$ whose continuation is novel. This fraction closely tracks the frequency dependence of the gap: the dashed curves in Figure~\ref{fig:frequency} are the directly measured $f$, rescaled and shifted,
\begin{equation}
\gap_t(F)\approx a_t f(F)+b_t,
\label{eq:coverage-main}
\end{equation}
with one scale and one offset per branch at training step $t$. Good--Turing estimation connects this novel fraction to singleton counts in the corpus; heavy-tailed conditional distributions provide a reason for its slow, approximately power-law decay. The detailed derivation and the terms beyond simple proportionality are in Appendix~\ref{app:math}. Thus the frequency trend points to a property of the text: rare contexts leave more plausible continuations unseen. It remains to explain why those continuations are harmed by replay.

\begin{figure}[!t]
\centering
\begin{subfigure}[b]{0.58\linewidth}
\centering
\includegraphics[width=\linewidth]{figures/walk_their_dog_loss.pdf}
\subcaption{}
\label{fig:walk-their-dog}
\end{subfigure}\hfill
\begin{subfigure}[b]{0.38\linewidth}
\centering
\includegraphics[width=\linewidth]{figures/fig09_crowd_ce_seen_novel.pdf}
\subcaption{}
\label{fig:crowd-ce}
\end{subfigure}
\caption{Crowding-out: (a) context-only continuation losses for ``walk their dog''; (b) validation cross-entropy of train-seen and novel continuations, pooled over contexts.}
\label{fig:crowding}
\end{figure}

\subsection{Crowding-out under replay}
\label{sec:sharpening}
A finite corpus does not fully represent the continuation distribution of each context: a context whose population distribution assigns equal probability to $C$ and $D$ may appear three times with $C$ and once with $D$ in training, so fitting the sample underpredicts $D$ on fresh data. The replay trajectory adds a further clue. After one complete pass, the corpus and its novel fraction are fixed, yet the gap continues to grow, by comparable amounts over the first few passes (Figure~\ref{fig:staircase}; Appendix~\ref{app:replay-clue}). The missing coverage alone does not specify this continuing damage. We hypothesize that each repetition reinforces the seen continuations and progressively suppresses the alternatives.

This would make forking ordinary overfitting confined to very small problems. For a context $c_i$, call the next-token logits computed from it its \emph{context subspace}, and split the vocabulary into the continuations of $c_i$ seen in training, those that appear only in validation, and the rest (Appendix~\ref{app:crowding-derivation}). If training shows a single continuation $C$, the loss in that subspace is $-\ln p_C$; an idealized step directly in logit coordinates raises the logit of $C$ by $\eta(1-p_C)$ and, through the softmax, lowers every other logit by $\eta p_v$. Mass flows from the unseen continuations onto the seen one, so the seen-to-unseen odds grow by a factor above one at every pass with nothing to stop them while $p_C<1$ (Appendix Equation~\eqref{eq:zt}). We call this \emph{crowding-out}; it provides an idealized account of the growth in the novel-continuation penalty with replay.

Figure~\ref{fig:walk-their-dog} shows the picture on one context. ``walk their dog'' is followed by \emph{down} once in training and, in the fixed validation window, once each by \emph{daily} and \emph{for}. Probed with context-only input, the probability of \emph{down} rises by more than an order of magnitude over five passes while \emph{daily} and \emph{for} are compressed---\emph{daily} by orders of magnitude---and so is the total mass outside the continuations observed in either sample. Figure~\ref{fig:crowd-ce} shows that the context is typical. Pooling validation positions whose bigram context occurs in training, split by whether the continuation was seen there (405{,}058 positions) or not (158{,}790), both models do well on seen continuations over the first three passes: the control's cross-entropy falls from 4.3 to 2.8 nats and the both-table model's stays near 3.4. On novel continuations they diverge: the control still improves, 6.9 to 5.6 nats, while the both-table model deteriorates from 7.5 to 9.0 across the first three epoch ends. This continued deterioration accompanies the gap staircase in Figure~\ref{fig:staircase}. Consistently, at epoch~3 the both-table model's predictive entropy on validation is 2.93 nats against a cross-entropy of 5.26, whereas the control's are 3.59 and 3.58 (same runs as Figure~\ref{fig:crowd-ce}): far more confident than accurate, the signature of a model that has committed to few continuations per context \citep{carlsson2025hyperfitting}. Crowding-out also explains why freezing the tables does not stop forking (Section~\ref{sec:interventions}): the subspace coordinate is already in place, and the backbone can continue the descent on its own.

This leaves one question: why does the same replay improve the no-gram model on novel continuations, while the $n$-gram model grows less accurate? The next experiment tests whether the table allows sharpening to proceed more independently within each context.

\subsection{Testing the locality of context subspaces}
\label{sec:context-locality}
\begin{figure}[!t]
\centering
\begin{subfigure}[b]{0.45\linewidth}
\includegraphics[width=\linewidth]{figures/fig07a_context_mask.pdf}
\subcaption{Gap contribution.}
\label{fig:ctxmask-gap}
\end{subfigure}\hfill
\begin{subfigure}[b]{0.45\linewidth}
\includegraphics[width=\linewidth]{figures/fig07b_seen_mass.pdf}
\subcaption{Mass on train-seen continuations.}
\label{fig:ctxmask-mass}
\end{subfigure}
\caption{Original context loss-masking diagnostic. The displayed trajectories motivate the locality test; the unresolved comparison with registered reruns is described in Appendix~\ref{app:mask-reconciliation}.}
\label{fig:context-subspace-mask}
\end{figure}
The mechanism of Section~\ref{sec:sharpening} suggests that the $n$-gram module gives each context a coordinate on which sharpening can accumulate. How independently do these coordinates evolve? Context-loss masking tests whether removing a context's own updates stops its gap growth while other contexts keep training. The original diagnostic holds out a random 2\% of training trigram contexts and excludes the next-token loss at their positions from the start of epoch~2 or epoch~3 onward, keeping inputs, trainable parameters and the loss denominator unchanged; a disjoint 2\% serves as unmasked peers, and each masked run is paired with normal training from the same checkpoint. Figure~\ref{fig:ctxmask-gap} plots their gap contribution $C_S=\sum_{i\in S}q_i[\bar\ell_{\mathrm{val},i}-\bar\ell_{\mathrm{tr},i}]$ on fixed batches. Until the mask fires the masked curves track normal training; afterwards they freeze---for two and three further passes---while the unmasked and normal curves take a fresh jump at every replay boundary. In these displayed trajectories, the gap of a masked context stops growing when its own replay stops, although the remaining 98\% of contexts keep training. This is the pattern expected if gap growth is largely local.

Figure~\ref{fig:ctxmask-mass} follows the sharpening coordinate itself, the predictive mass on the continuations of each context that training has seen. In the n-gram model the unmasked contexts pour mass into their seen continuations pass after pass ($0.43\to0.69$ by epoch~4) while the masked ones freeze ($0.43\to0.40$): the displayed sharpening largely stops when a context's own loss is removed. Without the $n$-gram module the mask changes nothing ($0.22\to0.39$ masked vs.\ $0.41$ unmasked), and replay raises the mass on seen and validation-only continuations alike ($0.013\to0.06$), consistent with learning through parameters shared across contexts. These trajectories suggest weakly interacting context subspaces. However, registered masking reruns did not reproduce the arrest of gap growth, and the original figure's protocol remains to be reconciled with them (Appendix~\ref{app:mask-reconciliation}). Context-specific sharpening is therefore a mechanism hypothesis; independence across contexts is not yet established.

